% Options for packages loaded elsewhere
\PassOptionsToPackage{unicode}{hyperref}
\PassOptionsToPackage{hyphens}{url}
%
\documentclass[
  12pt,
  a4paper,
  DIV=11,
  numbers=noendperiod]{scrreprt}

\usepackage{amsmath,amssymb}
\usepackage{setspace}
\usepackage{iftex}
\ifPDFTeX
  \usepackage[T1]{fontenc}
  \usepackage[utf8]{inputenc}
  \usepackage{textcomp} % provide euro and other symbols
\else % if luatex or xetex
  \usepackage{unicode-math}
  \defaultfontfeatures{Scale=MatchLowercase}
  \defaultfontfeatures[\rmfamily]{Ligatures=TeX,Scale=1}
\fi
\usepackage{lmodern}
\ifPDFTeX\else  
    % xetex/luatex font selection
\fi
% Use upquote if available, for straight quotes in verbatim environments
\IfFileExists{upquote.sty}{\usepackage{upquote}}{}
\IfFileExists{microtype.sty}{% use microtype if available
  \usepackage[]{microtype}
  \UseMicrotypeSet[protrusion]{basicmath} % disable protrusion for tt fonts
}{}
\usepackage{xcolor}
\setlength{\emergencystretch}{3em} % prevent overfull lines
\setcounter{secnumdepth}{5}
% Make \paragraph and \subparagraph free-standing
\makeatletter
\ifx\paragraph\undefined\else
  \let\oldparagraph\paragraph
  \renewcommand{\paragraph}{
    \@ifstar
      \xxxParagraphStar
      \xxxParagraphNoStar
  }
  \newcommand{\xxxParagraphStar}[1]{\oldparagraph*{#1}\mbox{}}
  \newcommand{\xxxParagraphNoStar}[1]{\oldparagraph{#1}\mbox{}}
\fi
\ifx\subparagraph\undefined\else
  \let\oldsubparagraph\subparagraph
  \renewcommand{\subparagraph}{
    \@ifstar
      \xxxSubParagraphStar
      \xxxSubParagraphNoStar
  }
  \newcommand{\xxxSubParagraphStar}[1]{\oldsubparagraph*{#1}\mbox{}}
  \newcommand{\xxxSubParagraphNoStar}[1]{\oldsubparagraph{#1}\mbox{}}
\fi
\makeatother


\providecommand{\tightlist}{%
  \setlength{\itemsep}{0pt}\setlength{\parskip}{0pt}}\usepackage{longtable,booktabs,array}
\usepackage{calc} % for calculating minipage widths
% Correct order of tables after \paragraph or \subparagraph
\usepackage{etoolbox}
\makeatletter
\patchcmd\longtable{\par}{\if@noskipsec\mbox{}\fi\par}{}{}
\makeatother
% Allow footnotes in longtable head/foot
\IfFileExists{footnotehyper.sty}{\usepackage{footnotehyper}}{\usepackage{footnote}}
\makesavenoteenv{longtable}
\usepackage{graphicx}
\makeatletter
\newsavebox\pandoc@box
\newcommand*\pandocbounded[1]{% scales image to fit in text height/width
  \sbox\pandoc@box{#1}%
  \Gscale@div\@tempa{\textheight}{\dimexpr\ht\pandoc@box+\dp\pandoc@box\relax}%
  \Gscale@div\@tempb{\linewidth}{\wd\pandoc@box}%
  \ifdim\@tempb\p@<\@tempa\p@\let\@tempa\@tempb\fi% select the smaller of both
  \ifdim\@tempa\p@<\p@\scalebox{\@tempa}{\usebox\pandoc@box}%
  \else\usebox{\pandoc@box}%
  \fi%
}
% Set default figure placement to htbp
\def\fps@figure{htbp}
\makeatother

\usepackage{indentfirst}
\usepackage{float}
\floatplacement{figure}{H}
\IfFileExists{plex-otf.sty}{
  %% Full TeXlive
  \usepackage[
    % math,
    RM={Scale=0.94},SS={Scale=0.94},SScon={Scale=0.94},TT={Scale=MatchLowercase,FakeStretch=0.9},
    DefaultFeatures={Ligatures=Common}
  ]{plex-otf}
  % \usepackage[Scale=MatchUppercase]{juliamono}
}{
  %% TinyTeX
  \usepackage{libertine}
}
\KOMAoption{captions}{tableheading}
\makeatletter
\@ifpackageloaded{caption}{}{\usepackage{caption}}
\AtBeginDocument{%
\ifdefined\contentsname
  \renewcommand*\contentsname{Содержание}
\else
  \newcommand\contentsname{Содержание}
\fi
\ifdefined\listfigurename
  \renewcommand*\listfigurename{Список иллюстраций}
\else
  \newcommand\listfigurename{Список иллюстраций}
\fi
\ifdefined\listtablename
  \renewcommand*\listtablename{Список таблиц}
\else
  \newcommand\listtablename{Список таблиц}
\fi
\ifdefined\figurename
  \renewcommand*\figurename{Рисунок}
\else
  \newcommand\figurename{Рисунок}
\fi
\ifdefined\tablename
  \renewcommand*\tablename{Таблица}
\else
  \newcommand\tablename{Таблица}
\fi
}
\@ifpackageloaded{float}{}{\usepackage{float}}
\floatstyle{ruled}
\@ifundefined{c@chapter}{\newfloat{codelisting}{h}{lop}}{\newfloat{codelisting}{h}{lop}[chapter]}
\floatname{codelisting}{Список}
\newcommand*\listoflistings{\listof{codelisting}{Листинги}}
\makeatother
\makeatletter
\makeatother
\makeatletter
\@ifpackageloaded{caption}{}{\usepackage{caption}}
\@ifpackageloaded{subcaption}{}{\usepackage{subcaption}}
\makeatother

\ifLuaTeX
\usepackage[bidi=basic]{babel}
\else
\usepackage[bidi=default]{babel}
\fi
\babelprovide[main,import]{russian}
\babelprovide[import]{english}
% get rid of language-specific shorthands (see #6817):
\let\LanguageShortHands\languageshorthands
\def\languageshorthands#1{}
\usepackage[style=gost-numeric,backend=biber,langhook=extras,autolang=other*]{biblatex}
\addbibresource{bib/cite.bib}
\usepackage{csquotes}
\usepackage{bookmark}

\IfFileExists{xurl.sty}{\usepackage{xurl}}{} % add URL line breaks if available
\urlstyle{same} % disable monospaced font for URLs
\hypersetup{
  pdftitle={Отчёт по лабораторной работе №6},
  pdfauthor={Пакавира Арсениу Висенте Луиш},
  pdflang={ru-RU},
  hidelinks,
  pdfcreator={LaTeX via pandoc}}


\title{Отчёт по лабораторной работе №6}
\usepackage{etoolbox}
\makeatletter
\providecommand{\subtitle}[1]{% add subtitle to \maketitle
  \apptocmd{\@title}{\par {\large #1 \par}}{}{}
}
\makeatother
\subtitle{Дисциплина: Администрирование сетевых подсистем}
\author{Пакавира Арсениу Висенте Луиш}
\date{}

\begin{document}
\maketitle

\renewcommand*\contentsname{Содержание}
{
\setcounter{tocdepth}{1}
\tableofcontents
}
\listoffigures
\listoftables

\setstretch{1.5}
\chapter{Цель
работы}\label{ux446ux435ux43bux44c-ux440ux430ux431ux43eux442ux44b}

Приобретение практических навыков по установке и конфигурированию
системы управления базами данных на примере программного обеспечения
MariaDB.

\chapter{Выполнение лабораторной
работы}\label{ux432ux44bux43fux43eux43bux43dux435ux43dux438ux435-ux43bux430ux431ux43eux440ux430ux442ux43eux440ux43dux43eux439-ux440ux430ux431ux43eux442ux44b}

\section{Установка
MariaDB}\label{ux443ux441ux442ux430ux43dux43eux432ux43aux430-mariadb}

На виртуальной машине \texttt{server} выполнен переход в режим
суперпользователя и произведена установка пакетов \texttt{mariadb} и
\texttt{mariadb-server} с использованием менеджера пакетов \texttt{dnf}
(\textbf{?@fig-1}). В процессе установки были автоматически подключены
зависимости и выполнена загрузка необходимых компонентов.

\includegraphics[width=0.7\linewidth,height=\textheight,keepaspectratio]{image/1.png}
\href{https://github.com/ArsenioPacavira99/study_2025-2026_net-os-admin-12/tree/master/labs/lab06/report/image/1.png}{Нажмите
здесь, чтобы открыть ссылку.}

После установки выполнен запуск службы \texttt{mariadb} и добавление её
в автозагрузку с помощью \texttt{systemctl\ start\ mariadb} и
\texttt{systemctl\ enable\ mariadb} (\textbf{?@fig-2}). Созданы
символические ссылки в каталогах systemd, что подтверждает успешную
регистрацию службы.

\includegraphics[width=0.7\linewidth,height=\textheight,keepaspectratio]{image/2.png}
\href{https://github.com/ArsenioPacavira99/study_2025-2026_net-os-admin-12/tree/master/labs/lab06/report/image/2.png}{Нажмите
здесь, чтобы открыть ссылку.}

Далее проверено, что процесс \texttt{mariadb} прослушивает порт 3306, с
использованием команды \texttt{ss\ -tulpen\ \textbar{}\ grep\ 3306}
(\textbf{?@fig-2}). В выводе отображён процесс \texttt{mariadb},
находящийся в состоянии LISTEN на порту 3306, что подтверждает
корректный запуск сервера базы данных.

\includegraphics[width=0.7\linewidth,height=\textheight,keepaspectratio]{image/2.png}
\href{https://github.com/ArsenioPacavira99/study_2025-2026_net-os-admin-12/tree/master/labs/lab06/report/image/2.png}{Нажмите
здесь, чтобы открыть ссылку.}

Выполнен скрипт первичной настройки безопасности
\texttt{mysql\_secure\_installation}. В ходе диалога включена
аутентификация через \texttt{unix\_socket}, установлен пароль для
пользователя \texttt{root} СУБД, удалены анонимные пользователи
(\textbf{?@fig-3}), запрещён удалённый доступ для root и удалена
тестовая база данных. После применения изменений произведена
перезагрузка таблиц привилегий. (\textbf{?@fig-3-1})

\includegraphics[width=0.7\linewidth,height=\textheight,keepaspectratio]{image/3.png}
\href{https://github.com/ArsenioPacavira99/study_2025-2026_net-os-admin-12/tree/master/labs/lab06/report/image/3.png}{Нажмите
здесь, чтобы открыть ссылку.}

\includegraphics[width=0.7\linewidth,height=\textheight,keepaspectratio]{image/3-1.png}
\href{https://github.com/ArsenioPacavira99/study_2025-2026_net-os-admin-12/tree/master/labs/lab06/report/image/3-1.png}{Нажмите
здесь, чтобы открыть ссылку.}

После завершения настройки выполнен вход в систему управления базами
данных от имени администратора с помощью \texttt{mysql\ -u\ root\ -p}
(\textbf{?@fig-4}). Отображена информация о версии сервера MariaDB и
параметрах подключения.

\includegraphics[width=0.7\linewidth,height=\textheight,keepaspectratio]{image/4.png}
\href{https://github.com/ArsenioPacavira99/study_2025-2026_net-os-admin-12/tree/master/labs/lab06/report/image/4.png}{Нажмите
здесь, чтобы открыть ссылку.}

В интерактивной оболочке MariaDB выполнена команда для просмотра списка
доступных клиентских команд (\textbf{?@fig-5}). В выводе представлен
перечень управляющих команд консольного клиента.

\includegraphics[width=0.7\linewidth,height=\textheight,keepaspectratio]{image/5.png}
\href{https://github.com/ArsenioPacavira99/study_2025-2026_net-os-admin-12/tree/master/labs/lab06/report/image/5.png}{Нажмите
здесь, чтобы открыть ссылку.}

С использованием SQL-запроса \texttt{SHOW\ DATABASES;} получен список
доступных системных баз данных (\textbf{?@fig-6}). В системе
присутствуют базы данных: \texttt{information\_schema}, \texttt{mysql},
\texttt{performance\_schema}.

\includegraphics[width=0.7\linewidth,height=\textheight,keepaspectratio]{image/6.png}
\href{https://github.com/ArsenioPacavira99/study_2025-2026_net-os-admin-12/tree/master/labs/lab06/report/image/6.png}{Нажмите
здесь, чтобы открыть ссылку.}

\section{Конфигурация кодировки
символов}\label{ux43aux43eux43dux444ux438ux433ux443ux440ux430ux446ux438ux44f-ux43aux43eux434ux438ux440ux43eux432ux43aux438-ux441ux438ux43cux432ux43eux43bux43eux432}

Выполнен вход в MariaDB от имени администратора базы данных с
использованием команды \texttt{mysql\ -u\ root\ -p}. После
аутентификации выведена информация о версии сервера и установленном
соединении (\textbf{?@fig-7}).

\includegraphics[width=0.7\linewidth,height=\textheight,keepaspectratio]{image/7.png}
\href{https://github.com/ArsenioPacavira99/study_2025-2026_net-os-admin-12/tree/master/labs/lab06/report/image/7.png}{Нажмите
здесь, чтобы открыть ссылку.}

Из интерактивной оболочки выполнена команда \texttt{status}, которая
отобразила текущее состояние сервера (\textbf{?@fig-7}). В выводе
указаны: идентификатор соединения (Connection id), текущий пользователь
(root@localhost), способ подключения (Localhost via UNIX socket), версия
сервера (10.5.29-MariaDB), а также параметры кодировки. До изменения
конфигурации значения \texttt{Server\ characterset} и
\texttt{Db\ characterset} установлены в \texttt{latin1}, а
\texttt{Client\ characterset} и \texttt{Conn.\ characterset} --- в
\texttt{utf8}.

Далее в каталоге \texttt{/etc/my.cnf.d} создан файл \texttt{utf8.cnf} и
открыт для редактирования (\textbf{?@fig-8}).

\includegraphics[width=0.7\linewidth,height=\textheight,keepaspectratio]{image/8.png}
\href{https://github.com/ArsenioPacavira99/study_2025-2026_net-os-admin-12/tree/master/labs/lab06/report/image/8.png}{Нажмите
здесь, чтобы открыть ссылку.}

В файл добавлена конфигурация:

\begin{verbatim}
[client]
default-character-set = utf8
[mysqld]
character-set-server = utf8
\end{verbatim}

Данные параметры задают кодировку \texttt{utf8} по умолчанию для
клиентских подключений и для сервера MariaDB (\textbf{?@fig-9}).

\includegraphics[width=0.7\linewidth,height=\textheight,keepaspectratio]{image/9.png}
\href{https://github.com/ArsenioPacavira99/study_2025-2026_net-os-admin-12/tree/master/labs/lab06/report/image/9.png}{Нажмите
здесь, чтобы открыть ссылку.}

После внесения изменений выполнен перезапуск службы MariaDB командой
\texttt{systemctl\ restart\ mariadb}, затем произведён повторный вход в
систему и выполнена команда \texttt{status} (\textbf{?@fig-10}).

\includegraphics[width=0.7\linewidth,height=\textheight,keepaspectratio]{image/10.png}
\href{https://github.com/ArsenioPacavira99/study_2025-2026_net-os-admin-12/tree/master/labs/lab06/report/image/10.png}{Нажмите
здесь, чтобы открыть ссылку.}

В результате изменения конфигурации параметры
\texttt{Server\ characterset} и \texttt{Db\ characterset} изменились с
\texttt{latin1} на \texttt{utf8}. Параметры
\texttt{Client\ characterset} и \texttt{Conn.\ characterset} также
установлены в \texttt{utf8}. Это подтверждает корректное применение
новой конфигурации кодировки после перезапуска сервера.

\section{Создание базы
данных}\label{ux441ux43eux437ux434ux430ux43dux438ux435-ux431ux430ux437ux44b-ux434ux430ux43dux43dux44bux445}

Выполнен вход в MariaDB под пользователем \texttt{root} с использованием
команды \texttt{mysql\ -u\ root\ -p} (\textbf{?@fig-11}). После
аутентификации отображена информация о версии сервера и установленном
соединении.

\includegraphics[width=0.7\linewidth,height=\textheight,keepaspectratio]{image/11.png}
\href{https://github.com/ArsenioPacavira99/study_2025-2026_net-os-admin-12/tree/master/labs/lab06/report/image/11.png}{Нажмите
здесь, чтобы открыть ссылку.}

Создана база данных \texttt{addressbook} с кодировкой \texttt{utf8} и
сопоставлением \texttt{utf8\_general\_ci}, затем выполнен переход к ней
и проверено отсутствие таблиц командой \texttt{SHOW\ TABLES;}
(\textbf{?@fig-12}). База данных успешно создана, на момент проверки
таблицы отсутствуют.

Создана таблица \texttt{city} с полями \texttt{name} и \texttt{city}
типа \texttt{VARCHAR(40)}. В таблицу добавлены записи: Иванов ---
Москва, Петров --- Сочи, Сидоров --- Дубна.

\includegraphics[width=0.7\linewidth,height=\textheight,keepaspectratio]{image/12.png}
\href{https://github.com/ArsenioPacavira99/study_2025-2026_net-os-admin-12/tree/master/labs/lab06/report/image/12.png}{Нажмите
здесь, чтобы открыть ссылку.}

Выполнен запрос \texttt{SELECT\ *\ FROM\ city;} (\textbf{?@fig-13}).
Результат запроса отображает три строки, соответствующие введённым
данным, что подтверждает корректность создания таблицы и вставки
записей.

\includegraphics[width=0.7\linewidth,height=\textheight,keepaspectratio]{image/13.png}
\href{https://github.com/ArsenioPacavira99/study_2025-2026_net-os-admin-12/tree/master/labs/lab06/report/image/13.png}{Нажмите
здесь, чтобы открыть ссылку.}

Создан пользователь \texttt{user@\textquotesingle{}\%\textquotesingle{}}
с паролем и предоставлены права \texttt{SELECT}, \texttt{INSERT},
\texttt{UPDATE}, \texttt{DELETE} на базу данных \texttt{addressbook}.
После выполнения \texttt{FLUSH\ PRIVILEGES;} изменения прав вступили в
силу. Командой \texttt{DESCRIBE\ city;} получена структура таблицы
(\textbf{?@fig-14}). Таблица содержит два поля: \texttt{name} и
\texttt{city}, оба типа \texttt{VARCHAR(40)}, допускающие значение
\texttt{NULL}.

\includegraphics[width=0.7\linewidth,height=\textheight,keepaspectratio]{image/14.png}
\href{https://github.com/ArsenioPacavira99/study_2025-2026_net-os-admin-12/tree/master/labs/lab06/report/image/14.png}{Нажмите
здесь, чтобы открыть ссылку.}

После выхода из MariaDB выполнена проверка списка баз данных с помощью
\texttt{mysqlshow\ -u\ root\ -p}, а также просмотр таблиц базы
\texttt{addressbook} под пользователями \texttt{root} и \texttt{bahis}
(\textbf{?@fig-15}). В списке присутствует база \texttt{addressbook}, а
в её составе отображается таблица \texttt{city}, что подтверждает
корректность создания базы данных и настройки прав доступа.

\includegraphics[width=0.7\linewidth,height=\textheight,keepaspectratio]{image/15.png}
\href{https://github.com/ArsenioPacavira99/study_2025-2026_net-os-admin-12/tree/master/labs/lab06/report/image/15.png}{Нажмите
здесь, чтобы открыть ссылку.}

\section{Резервные
копии}\label{ux440ux435ux437ux435ux440ux432ux43dux44bux435-ux43aux43eux43fux438ux438}

На виртуальной машине \texttt{server} создан каталог для хранения
резервных копий с использованием команды \texttt{mkdir\ -p\ /var/backup}
(\textbf{?@fig-16}). Каталог предназначен для размещения дампов базы
данных.

Выполнено создание полной резервной копии базы данных
\texttt{addressbook} с помощью команды
\texttt{mysqldump\ -u\ root\ -p\ addressbook\ \textgreater{}\ /var/backup/addressbook.sql}
(\textbf{?@fig-16}). В результате сформирован SQL-файл, содержащий
структуру таблиц и данные базы.

Создана сжатая резервная копия базы данных с использованием конвейера и
утилиты \texttt{gzip}:
\texttt{mysqldump\ -u\ root\ -p\ addressbook\ \textbar{}\ gzip\ \textgreater{}\ /var/backup/addressbook.sql.gz}
(\textbf{?@fig-16}). В результате получен архивированный файл дампа.

Дополнительно выполнено создание сжатой резервной копии с указанием даты
и времени создания файла:
\texttt{mysqldump\ -u\ root\ -p\ addressbook\ \textbar{}\ gzip\ \textgreater{}\ \$(date\ +/var/backup/addressbook.\%Y\%m\%d.\%H\%M\%S.sql.gz)}
(\textbf{?@fig-16}). Имя файла автоматически формируется с добавлением
временной метки.

Произведено восстановление базы данных из обычной резервной копии
командой
\texttt{mysql\ -u\ root\ -p\ addressbook\ \textless{}\ /var/backup/addressbook.sql}
(\textbf{?@fig-16}). Данные базы импортированы из SQL-файла.

Выполнено восстановление базы данных из сжатой резервной копии с
использованием \texttt{zcat}:
\texttt{zcat\ /var/backup/addressbook.sql.gz\ \textbar{}\ mysql\ -u\ root\ -p\ addressbook}
(\textbf{?@fig-16}). Архив предварительно распакован в поток вывода и
передан в клиент \texttt{mysql}, что подтверждает корректность процедуры
восстановления.

\includegraphics[width=0.7\linewidth,height=\textheight,keepaspectratio]{image/16.png}
\href{https://github.com/ArsenioPacavira99/study_2025-2026_net-os-admin-12/tree/master/labs/lab06/report/image/16.png}{Нажмите
здесь, чтобы открыть ссылку.}

\section{Внесение изменений в настройки внутреннего окружения
виртуальной
машины}\label{ux432ux43dux435ux441ux435ux43dux438ux435-ux438ux437ux43cux435ux43dux435ux43dux438ux439-ux432-ux43dux430ux441ux442ux440ux43eux439ux43aux438-ux432ux43dux443ux442ux440ux435ux43dux43dux435ux433ux43e-ux43eux43aux440ux443ux436ux435ux43dux438ux44f-ux432ux438ux440ux442ux443ux430ux43bux44cux43dux43eux439-ux43cux430ux448ux438ux43dux44b}

На виртуальной машине \texttt{server} выполнен переход в каталог
\texttt{/vagrant/provision/server}. Создан каталог \texttt{mysql} со
структурой подкаталогов \texttt{etc/my.cnf.d} и \texttt{var/backup}. В
созданные каталоги скопированы файл конфигурации \texttt{utf8.cnf} и
резервные копии базы данных \texttt{addressbook} (\textbf{?@fig-17}).

\includegraphics[width=0.7\linewidth,height=\textheight,keepaspectratio]{image/17.png}
\href{https://github.com/ArsenioPacavira99/study_2025-2026_net-os-admin-12/tree/master/labs/lab06/report/image/17.png}{Нажмите
здесь, чтобы открыть ссылку.}

В каталоге \texttt{/vagrant/provision/server} создан исполняемый файл
\texttt{mysql.sh}, которому назначены права на выполнение. В файл
добавлен скрипт автоматической установки и настройки MariaDB
(\textbf{?@fig-18}). Скрипт выполняет установку пакетов \texttt{mariadb}
и \texttt{mariadb-server}, копирование конфигурационных файлов, создание
каталога \texttt{/var/backup}, запуск и включение службы
\texttt{mariadb}, выполнение \texttt{mysql\_secure\_installation} в
неинтерактивном режиме, создание базы данных \texttt{addressbook} и
восстановление данных из резервной копии.

\includegraphics[width=0.7\linewidth,height=\textheight,keepaspectratio]{image/18.png}
\href{https://github.com/ArsenioPacavira99/study_2025-2026_net-os-admin-12/tree/master/labs/lab06/report/image/18.png}{Нажмите
здесь, чтобы открыть ссылку.}

В конфигурационный файл \texttt{Vagrantfile} добавлена секция
провижининга для сервера, обеспечивающая автоматическое выполнение
скрипта \texttt{mysql.sh} при запуске виртуальной машины
(\textbf{?@fig-19}). Запись использует тип \texttt{shell}, параметр
\texttt{preserve\_order:\ true} и путь к скрипту
\texttt{provision/server/mysql.sh}.

\includegraphics[width=0.7\linewidth,height=\textheight,keepaspectratio]{image/19.png}
\href{https://github.com/ArsenioPacavira99/study_2025-2026_net-os-admin-12/tree/master/labs/lab06/report/image/19.png}{Нажмите
здесь, чтобы открыть ссылку.}

\chapter{Выводы}\label{ux432ux44bux432ux43eux434ux44b}

В ходе работы выполнена установка и базовая настройка сервера баз данных
MariaDB на виртуальной машине \texttt{server}. Проверена корректность
запуска службы и прослушивания стандартного порта 3306.

Произведена настройка параметров безопасности с использованием
\texttt{mysql\_secure\_installation}: установлен пароль пользователя
\texttt{root} СУБД, удалены анонимные пользователи, запрещён удалённый
доступ администратора и удалена тестовая база данных.

Настроена кодировка \texttt{utf8} на уровне клиента и сервера путём
изменения конфигурационных файлов и перезапуска службы. Корректность
применения параметров подтверждена через команду \texttt{status}.

Создана база данных \texttt{addressbook}, таблица \texttt{city},
добавлены тестовые записи и проверена их выборка SQL-запросом. Настроен
пользователь с ограниченными правами доступа к базе данных.

Освоены процедуры резервного копирования и восстановления базы данных с
использованием \texttt{mysqldump}, включая создание сжатых архивов.

Автоматизирован процесс установки и настройки MariaDB посредством
скрипта \texttt{mysql.sh}, подключённого к механизму provisioning в
\texttt{Vagrantfile}, что обеспечивает воспроизводимость конфигурации
при повторном развёртывании виртуальной машины.

\chapter{Ответы на контрольные
вопросы}\label{ux43eux442ux432ux435ux442ux44b-ux43dux430-ux43aux43eux43dux442ux440ux43eux43bux44cux43dux44bux435-ux432ux43eux43fux440ux43eux441ux44b}

\begin{enumerate}
\def\labelenumi{\arabic{enumi}.}
\tightlist
\item
  Какая команда отвечает за настройки безопасности в MariaDB?
\end{enumerate}

\begin{itemize}
\tightlist
\item
  Настройки безопасности в MariaDB обычно управляются с помощью команды
  mysql\_secure\_installation. Эта команда выполняет несколько шагов,
  включая установку пароля для пользователя root, удаление анонимных
  учетных записей, отключение удаленного входа для пользователя root и
  удаление тестовых баз данных.
\end{itemize}

\begin{enumerate}
\def\labelenumi{\arabic{enumi}.}
\setcounter{enumi}{1}
\tightlist
\item
  Как настроить MariaDB для доступа через сеть?
\end{enumerate}

\begin{itemize}
\tightlist
\item
  Для настройки MariaDB для доступа через сеть, вы можете
  отредактировать файл конфигурации MariaDB (обычно называемый my.cnf) и
  убедиться, что параметр bind-address установлен на IP-адрес, доступный
  в вашей сети. Также, убедитесь, что пользователь имеет права доступа
  извне, например, с использованием команды GRANT.
\end{itemize}

\begin{enumerate}
\def\labelenumi{\arabic{enumi}.}
\setcounter{enumi}{2}
\tightlist
\item
  Какая команда позволяет получить обзор доступных баз данных после
  входа в среду оболочки MariaDB?
\end{enumerate}

\begin{itemize}
\tightlist
\item
  SHOW DATABASES;
\end{itemize}

\begin{enumerate}
\def\labelenumi{\arabic{enumi}.}
\setcounter{enumi}{3}
\tightlist
\item
  Какая команда позволяет узнать, какие таблицы доступны в базе данных?
\end{enumerate}

\begin{itemize}
\tightlist
\item
  SHOW TABLES;
\end{itemize}

\begin{enumerate}
\def\labelenumi{\arabic{enumi}.}
\setcounter{enumi}{4}
\tightlist
\item
  Какая команда позволяет узнать, какие поля доступны в таблице? -
\end{enumerate}

\begin{itemize}
\tightlist
\item
  DESCRIBE table\_name;
\end{itemize}

\begin{enumerate}
\def\labelenumi{\arabic{enumi}.}
\setcounter{enumi}{5}
\tightlist
\item
  Какая команда позволяет узнать, какие записи доступны в таблице?
\end{enumerate}

\begin{itemize}
\tightlist
\item
  SELECT * FROM table\_name;
\end{itemize}

\begin{enumerate}
\def\labelenumi{\arabic{enumi}.}
\setcounter{enumi}{6}
\tightlist
\item
  Как удалить запись из таблицы?
\end{enumerate}

\begin{itemize}
\tightlist
\item
  DELETE FROM table\_name WHERE condition;, где condition - условие,
  определяющее, какие записи следует удалить.
\end{itemize}

\begin{enumerate}
\def\labelenumi{\arabic{enumi}.}
\setcounter{enumi}{7}
\tightlist
\item
  Где расположены файлы конфигурации MariaDB? Что можно настроить с их
  помощью?
\end{enumerate}

\begin{itemize}
\tightlist
\item
  Файлы конфигурации MariaDB обычно располагаются в различных местах в
  зависимости от системы, но основной файл - my.cnf. Он может быть в
  /etc/my.cnf, /etc/mysql/my.cnf или /usr/etc/my.cnf. С помощью этих
  файлов можно настроить различные параметры, такие как порт, пути к
  файлам данных, параметры безопасности и другие.
\end{itemize}

\begin{enumerate}
\def\labelenumi{\arabic{enumi}.}
\setcounter{enumi}{8}
\tightlist
\item
  Где располагаются файлы с базами данных MariaDB?
\end{enumerate}

\begin{itemize}
\tightlist
\item
  Файлы с базами данных MariaDB располагаются в директории данных.
  Обычно это /var/lib/mysql/ на Linux-системах.
\end{itemize}

\begin{enumerate}
\def\labelenumi{\arabic{enumi}.}
\setcounter{enumi}{9}
\tightlist
\item
  Как сделать резервную копию базы данных и затем её восстановить?
\end{enumerate}

\begin{itemize}
\tightlist
\item
  Для создания резервной копии базы данных используйте команду
  mysqldump. Например, mysqldump -u username -p dbname \textgreater{}
  backup.sql. Для восстановления базы данных из резервной копии
  используйте команду mysql -u username -p dbname \textless{}
  backup.sql.
\end{itemize}


\printbibliography



\end{document}
